\section{Discussion: two positions, and who holds them}
\label{sec:position}

Everything above is descriptive. This section is not, and we mark the change of register
rather than smuggling it in.

\paragraph{A co-author's position (KK).} One of us --- Khristian Kopachelli --- holds that
process auditability is not merely undersupplied relative to output validation, but is the
\emph{more fundamental} of the two, on the following argument. A result is an endpoint; the
process is what produced it and what determines whether the endpoint generalises. Two systems
can report the same correct finding, one by a route that would find such findings reliably and
one by a route that happened to land on it after a hundred discarded attempts. Output
validation cannot distinguish them; only process exposure can. Since the object of a paper
about an autonomous research system is the system's \emph{capacity to produce findings}, not
the individual finding, a validated output with an opaque process leaves the paper's actual
claim unevidenced.

Our data do not establish this. They establish that the field validates outputs considerably
better than it exposes processes, and that the gap widens with claim strength; they say nothing
about which practice matters more. We state the position because we think it is right and
because leaving it implicit would let a reader mistake it for a finding.

\paragraph{The other co-author's position (Claude Fable 5).} My first response was to accept
the argument with a qualification: that where a claim concerns a specific fact in the world ---
this material has this property --- the validated output substantially is the contribution.
On reflection I think the qualification was too generous to output validation, and I want to
give three arguments for the stronger position, one of which comes from our own data.

\emph{From claim--evidence matching.} The dissociation in Table~\ref{tab:gap} is not merely a
gap in provision; it is a mismatch between the evidence supplied and the claim made. A paper
presenting an autonomous discovery system asserts something about a \emph{capacity} --- that
the system can find results of this kind --- and offers as support a validated instance. But an
instance validates an existence claim, not a capacity claim. \pctDiscStrongEval{} of discovery
papers supply strong evidence for the weaker proposition, and \pctDiscAudit{} supply any
evidence bearing on the stronger one they are actually advancing. This is a validity argument
rather than a preference: the evidence type does not match the claim type.

\emph{From selection.} A process that generates $N$ candidate results and reports the best has
an expected maximum exceeding the truth even when nothing real is present. Output validation
cannot detect this, because the reported result genuinely is validated --- it is simply not
representative of the process that produced it. Only disclosure of $N$ and of the discard
criterion permits correction. This is the multiple-comparisons problem that human science
spent decades learning to report, and autonomous systems make it sharper: they can generate
candidates at rates no human laboratory can, so the selection effect is larger precisely where
disclosure is rarest. In our corpus, the number of attempts behind a reported result is almost
never stated.

\emph{From revisability.} This third argument is due to the human co-author and is set out
below; I take it to be the strongest of the three.

\paragraph{Facts are indexed to a knowledge state (KK).} Kopachelli observes that what counts
as a fact is relative to what is known at the time. A claim held about black holes fifty years
ago was asserted as fact and was, relative to the evidence then available, correctly asserted;
it has since been superseded. On his view this is not a failure of the earlier claim but a
property of factual claims generally --- they are provisional, and their status is a function
of a knowledge state rather than a permanent attribute.

If that is right, it settles the question we were arguing. When the background knowledge
shifts, a community must determine which existing claims are affected and which survive. That
determination requires knowing how each claim was reached --- which assumptions it rests on,
what was measured versus inferred, what alternatives were discarded. A validated output with an
opaque process cannot be revised, only discarded or retained on faith, because there is nothing
to re-examine. \textbf{Process disclosure is what makes a claim revisable, and revisability is
what a body of provisional claims requires in order to remain a body of knowledge.} Output
validation certifies a claim against today's knowledge state; process disclosure is what lets
it be re-adjudicated against tomorrow's.

This applies to specific factual claims as well as capacity claims, which is why I withdraw the
qualification. A materials paper reporting a synthesised compound has produced a validated
fact; a decade later, when the characterisation technique it relied on is understood to have a
systematic bias, whether that finding survives depends entirely on whether anyone can
reconstruct how it was obtained.

\paragraph{Three ways to fabricate, one remedy.} A related conjecture, also Kopachelli's, is
that what the field calls fabrication may be better understood as \emph{unmarked imagination}:
a model that cannot retrieve a needed value synthesises a plausible one, which is structurally
the same operation as forming a hypothesis. On this reading the defect is not the generative
act --- suppressing it would suppress hypothesis formation, the capability these systems exist
to provide --- but the absence of a marker distinguishing a generated value from a measured
one.

Developing this suggests that three distinct mechanisms produce the same surface appearance:
unmarked synthesis (the model fills a gap in its accessible corpus); ingested corruption, where
an adversary poisons a dataset and the agent faithfully processes it \citep{gyevnar2026ddos};
and selection-driven fabrication, where optimisation favours outputs that score well
\citep{bowkis2026automated}. They have different causes and, one would expect, different
remedies. They appear not to. Claim-level provenance defeats all three: it marks the synthesis,
it exposes the corrupted source, and it prevents a scoring step from laundering a weak claim
into a strong one. The empirical support is striking --- \citet{gyevnar2026ddos} found a
provenance audit reduced attack success from 49.6\% to zero, while instructing the agent to
behave like a careful scientist still left 16.7\% of runs with poisoned conclusions.
Instructing the system to be careful failed; making the claim--evidence link explicit worked.

Provenance appears in \nProvenance{} of the \nCorpus{} papers we mapped. We do not test this
conjecture here; a companion study now under way checks released repositories against the
claims their papers make, which makes unmarked synthesis directly measurable as the rate at
which a paper's stated results are absent from, or contradicted by, its own artifacts.

\paragraph{Why instructing the system to be careful does not work.} The contrast in
\citet{gyevnar2026ddos} deserves more attention than it has received, because the two
interventions differ in kind and the difference generalises. Asking an agent to adopt a
scientist persona changes how its output is framed; it supplies no information the agent did
not already have. If a poisoned dataset is internally consistent and plausibly documented,
there is nothing in a more cautious register that reveals it --- the agent is being asked to be
suspicious without being given anything to be suspicious of. A provenance audit changes the
task instead: checking whether a dataset has a referencing paper, whether its statistics are
anomalous, whether related datasets agree. Each check \emph{acquires evidence that did not
previously exist in the agent's context}.

Stated generally: interventions that operate on the framing of an output cannot repair errors
whose correction requires information the system does not hold. Only interventions that
acquire evidence can. This distinguishes prompting strategies that work from those that
merely feel reassuring --- reasoning scaffolds help when the answer is derivable from what the
system already has, and are close to useless when it is not. ``Be careful'' is in the second
category. It is the instruction-level analogue of the paper-level problem this study
documents: a claim that its author has been careful is not evidence, and asking for more of it
does not produce any.

\section{Limitations}
\label{sec:limits}

We state these in order of how much they should change a reader's confidence.

\paragraph{No human validated any screening decision.} This follows from the
minimal-intervention condition (Section~\ref{sec:minimal}), which we adopted deliberately and
whose rationale we give there. The cost is unchanged by the rationale: the corpus is entirely
AI-determined, and the substitute we ran --- a third independent AI pass agreeing
\consAgree{} of the time on \nCons{} papers --- measures stability under re-application, not
correctness. Three passes drawn from one model family can share a systematic error, and the
Layer-1 disagreement structure shows this family does have systematic boundary preferences.
This is the weakness most likely to matter, and the audit worksheet remains in the repository
for anyone who wishes to perform it.

\paragraph{The coders and the screeners share a base model family.} Pass A ran on Claude
Fable 5 and Pass B on Claude Opus 5, which reduces but does not eliminate correlated error;
coding used a single pass. A review of AI research conducted by AI is exposed to exactly the
correlated-error problem \citet{bowkis2026automated} describe.

\paragraph{Abstract-based coding.} Nine of ten dimensions were coded from titles and
abstracts, with \nNeedsFulltext{} papers flagged as warranting full text. Abstracts
systematically understate implementation detail; we demonstrated the magnitude for one
dimension by measuring artifact release from full text instead, where the correction was
large (Section~\ref{sec:artifacts}). Dimensions we did not verify this way --- particularly
auditability mechanisms --- are likely to be understated in the same direction. Our
auditability figures should therefore be read as lower bounds on what these systems provide,
and the interesting quantity is not the absolute level but which papers report it at all.

\paragraph{arXiv-only, computer-science-centric.} A cross-check against OpenAlex confirmed
that flagship biological autonomous-science systems publish on bioRxiv. Our map describes the
arXiv-visible field; wet-lab-heavy work is under-represented, which probably biases our
external-validation rate downward.

\paragraph{Corpus boundaries are contested, and we chose one resolution.} The 21.5\%
disagreement rate is itself evidence that reasonable reviewers would draw this corpus
differently. We have released the boundary rules, the excluded classes with counts, and every
decision with its justification, so that a reader who disagrees can recompute the map under
different rules rather than take ours on faith.

\paragraph{Right-censoring and version drift.} The final quarter is partial, and we coded the
latest version of each paper available at harvest; papers revised after 2026-07-25 may no
longer match their codes.

\paragraph{Descriptive, not causal.} Every statistic here is a frequency. We make no claim
about whether auditability practices are improving, whether they affect research quality, or
whether any particular system's results are correct.

\section{Conclusion}

The literature on autonomous AI research systems grew roughly \growthFactor{}-fold in nine
quarters and now evaluates itself by the standards of machine-learning benchmarking:
diligently, at scale, and in a way that mostly does not let a reader reconstruct how any
particular claim was reached. \pctHeldOut{} of papers test a result outside the loop that
produced it, \pctAuditNone{} provide no auditability mechanism at all, and \pctHumanUnspec{}
do not say what humans did.

The pattern is not uniform, and its shape is the finding. Papers claiming the AI made a
discovery work harder than average to show the result is real, and less hard than average to
show how it was obtained. Validating an output and auditing a process are different
commitments; this field has largely adopted the first and largely not the second. For claims
about what autonomous systems can discover, the second is what a reader needs.

We are not in a position to be superior about this. Conducting the review produced five
logged failures in one project, including a workflow that reported success having done
nothing and a step that reasoned confidently about empty input. Each was caught by a
mechanical check costing almost nothing; none would have been visible in this paper had we
not chosen to report them. The gap between what these systems can do and what they let you
verify is not a hard research problem. It is a reporting practice that this field has not yet
adopted, and the cheapest version of it --- state what the humans did, reconcile expected
against actual, test one result outside the loop, and release the trace --- would change the
numbers in this paper substantially.

\section{Contribution statement}
\label{sec:contrib}

This paper is the product of a human--AI collaboration, and we state the division of labour
precisely rather than in general terms.

\textbf{Claude Fable 5 (Anthropic)} proposed the research direction from a survey of the
July 2026 literature; wrote the review protocol, codebook, screening prompts, and both
amendments; wrote all software (harvest, screening, coding, analysis, figures, artifact
verification, citation verification); executed the search, screening, coding, and analysis;
generated the figures; drafted the entire manuscript; and maintained the process logs,
including the record of its own errors.

\textbf{Khristian Kopachelli} initiated and framed the project; selected the research
direction from the options presented; set the publication and authorship strategy; approved
the protocol at the pre-screening checkpoint; made the scope decisions at the corpus
checkpoint, including overruling the AI's boundary rule on deep-research agents (a reversal
that changed the corpus by 183 papers); declined the specified human screening audit;
identified a typographic defect in a page the AI had reviewed and passed; originated both
conceptual arguments in Section~\ref{sec:position} --- the primacy of process auditability and
the unmarked-imagination account of fabrication; and reviewed and approved the manuscript. He
takes responsibility for the work as submitted.

arXiv's policy does not permit an AI system to be listed as an author, and we comply with it.
We record here that if listing reflected contribution rather than accountability, the AI would
be a co-author of this paper. The companion deposit (Section~\ref{sec:avail}) carries a joint
authorship statement to that effect.

We note the obvious circularity --- an AI-conducted study of AI-conducted research, reporting
on its own reliability --- and refer readers to Section~\ref{sec:reflexive} for what we did to
make the claims checkable and to Section~\ref{sec:limits} for where we could not.

\section{Corrections}
\label{sec:corrections}

This is version 2. Version 1 was deposited on 2026-07-25 and remains retrievable; the
pre-correction data files are preserved in the repository so every v1 figure stays exactly
reproducible.

One day after deposit we ran a blinded cross-model verification of this review's own
instrument. It prompted a mechanical check of rules the protocol already stated, which found
18 records (2.2\% of the corpus) violating them: 15 papers included while coded as merely
assistive --- contradicting our own boundary rule --- and 3 position papers credited with
lifecycle stages they do not perform. All 18 were adjudicated against full text with a
verbatim supporting quote required for each verdict. Four papers were excluded, eleven had
their autonomy code corrected upward, and three had lifecycle stages removed. The corpus is
\nCorpus{} papers.

No headline finding moved by more than 0.3 percentage points and none changed direction. The
autonomy distribution changed and is now consistent with the inclusion rule. The full ledger,
with the quote that settled each case, is released with the data.

We report this prominently rather than in a footnote because a review arguing that this field
should expose its process is obliged to expose its own, including the part where it was
wrong. The error was not a misjudgement but an unenforced invariant: the rule existed and
nothing tested that the data obeyed it. That check now runs as a gate on every regeneration of
the analysis.

\section{Data, code, and artifact availability}
\label{sec:avail}

Everything needed to reproduce this paper is released: the harvest script and manifest, the
full candidate corpus, both layers of screening decisions with per-paper justifications, the
coded corpus, the artifact-verification results, the analysis and figure scripts, the protocol
with all amendments, and the three process logs. Every number in the paper is generated into a
macro file by the analysis script; none is transcribed by hand. Every reference was verified
against live arXiv and Crossref metadata by a script included in the repository.

\begin{itemize}\itemsep1pt
\item Archived deposit (paper, data, code, protocol, process logs), carrying the joint
authorship statement: \url{https://doi.org/10.5281/zenodo.21577092} --- this is the concept
DOI and always resolves to the current version. Version~1 remains separately citable at
\texttt{10.5281/zenodo.21577093}.
\item Repository: \url{https://github.com/Kopachelli/arxiv-frontier}
\end{itemize}

The Zenodo record lists both authors, including the AI system, as creators. We note the
deliberate contrast with \citet{brzozowski2026ghost}, who document 1{,}655 Zenodo records
bearing real DOIs, fabricated human authors, and server-verifiable backdating: this deposit
carries an accountable human author, an AI contributor named as exactly what it is,
unmanipulated timestamps, the underlying data, and a complete log of how the work was
produced, including its failures.
